% !TeX program = xelatex
\documentclass[UTF8,a4paper,zihao=-4,fontset=fandol]{ctexart}
\usepackage[margin=18mm,top=18mm,bottom=18mm,headheight=15pt,headsep=6mm]{geometry}
\usepackage{amsmath,booktabs,tabularx,array,multirow}
\usepackage[table]{xcolor}
\usepackage{tikz}
\usetikzlibrary{arrows.meta,positioning,calc,shapes.geometric,fit}
\usepackage{enumitem,fancyhdr,caption,float,hyperref}
\definecolor{Ink}{HTML}{18354A}
\definecolor{Teal}{HTML}{117F83}
\definecolor{Pale}{HTML}{EEF5F7}
\definecolor{Mist}{HTML}{F5F7F8}
\definecolor{Amber}{HTML}{986421}
\definecolor{Warm}{HTML}{FFF5E6}
\definecolor{Muted}{HTML}{5B6B76}
\definecolor{Line}{HTML}{CAD9DF}
\hypersetup{colorlinks=true,linkcolor=Ink,urlcolor=Teal,citecolor=Teal,
 pdftitle={ACG720赛题三配件与软件采购指南},pdfauthor={项目采购规划},bookmarksnumbered=true}
\urlstyle{same}
\ctexset{section={format=\Large\bfseries\color{Ink},beforeskip=0pt,afterskip=10pt},
 subsection={format=\large\bfseries\color{Ink},beforeskip=11pt,afterskip=6pt}}
\setlength{\parindent}{0pt}
\setlength{\parskip}{5pt}
\linespread{1.14}
\setlist[itemize]{leftmargin=1.55em,itemsep=2pt,topsep=3pt,parsep=0pt}
\setlength{\tabcolsep}{5pt}
\renewcommand{\arraystretch}{1.28}
\newcolumntype{Y}{>{\raggedright\arraybackslash}X}
\newcolumntype{L}[1]{>{\raggedright\arraybackslash}p{#1}}
\captionsetup{font=small,labelfont={bf,color=Ink},skip=5pt}
\pagestyle{fancy}\fancyhf{}
\fancyhead[L]{\small\color{Muted}ACG720 / 赛题三配件与软件}
\fancyhead[R]{\small\color{Muted}采购前核对版}
\fancyfoot[L]{\footnotesize\color{Muted}型号建议不等于板级兼容性承诺}
\fancyfoot[R]{\small\color{Ink}\thepage}
\renewcommand{\headrulewidth}{0.3pt}
\newcommand{\note}[2]{\par\smallskip\noindent\colorbox{Pale}{\parbox{\dimexpr\linewidth-2\fboxsep\relax}{\small\textbf{\color{Teal}#1}\quad #2}}\par\smallskip}
\newcommand{\warn}[2]{\par\smallskip\noindent\colorbox{Warm}{\parbox{\dimexpr\linewidth-2\fboxsep\relax}{\small\textbf{\color{Amber}#1}\quad #2}}\par\smallskip}
\newcommand{\newpagex}{\clearpage}
\newcommand{\tagline}[1]{{\small\color{Muted}#1}\par}
\tikzset{box/.style={draw=Line,fill=Pale,rounded corners=2pt,align=center,minimum height=10mm,inner sep=5pt,font=\small},
 main/.style={box,fill=Ink,text=white,draw=Ink},
 teal/.style={box,fill=Teal!10,draw=Teal!60},
 wait/.style={box,fill=Warm,draw=Amber!70,dashed},
 flow/.style={-{Latex[length=2mm]},thick,draw=Ink},
 opt/.style={flow,dashed,draw=Amber},
 label/.style={font=\footnotesize,text=Muted,align=center}}

\title{\textbf{\color{Ink}ACG720 赛题三配件与软件采购指南}\\[7pt]
 {\large FPGA 实时控制 + 树莓派拓展 AI 方案}}
\author{配件清单 · 分批采购 · 供电与接口 · 软件许可}
\date{依据 2026 高云 FPGA 赛道指南及 ACG720/GW5AT-138K 配套资料\\资料核验：2026 年 9 月 27 日}

\begin{document}
\maketitle
\thispagestyle{empty}
\vspace{-5mm}
\note{一句话结论}{先购买能复用现成例程的摄像头、基础传感器和树莓派开发附件；底盘、电机、驱动器与电池作为同一批次定型。树莓派只承担任务5、6的拓展视觉，不替代 FPGA 的实时感知、闭环控制和安全停车。}

\begin{figure}[H]\centering
\begin{tikzpicture}[x=1cm,y=1cm]
\node[main,text width=4cm] (root) at (0,0) {机器人系统};
\node[teal,text width=4cm] (fpga) at (-4,-1.5) {ACG720\;基础实时层};
\node[wait,text width=4cm] (pi) at (4,-1.5) {Raspberry Pi 5\;拓展 AI 层};
\node[box,text width=4cm,align=left] (f) at (-4,-3.8) {ACM5640-V5\\IMU + 4路 ToF\\编码器 + 电机驱动\\DDR3 + HDMI/TFT};
\node[box,text width=4cm,align=left] (p) at (4,-3.8) {Camera Module 3\\目标检测/跟踪\\颜色属性与搜索\\高层任务请求};
\draw[flow] (root)--(fpga);\draw[opt] (root)--(pi);
\draw[flow] (fpga)--(f);\draw[opt] (pi)--(p);
\draw[<->,dashed,thick,draw=Amber] (fpga)--node[above,label]{UART 起步\\以太网传图后加}(pi);
\end{tikzpicture}
\caption{推荐职责划分。安全闭环必须在树莓派失联时仍然有效。}\label{fig:architecture}
\end{figure}

\begin{tabularx}{\linewidth}{@{}L{2.6cm}Y r@{}}\toprule
\textbf{阅读位置} & \textbf{解决的问题} & \textbf{页码}\\\midrule
1 合规边界 & 树莓派能否使用，怎样使用风险最低 & \pageref{sec:rules}\\
2 首批采购 & 哪些型号现在可以买 & \pageref{sec:first}\\
3 运动与电源 & 哪些必须成套定型 & \pageref{sec:motion}\\
4 可选升级 & 什么暂时不要买 & \pageref{sec:upgrade}\\
5 软件与许可 & 哪些要注册、申请或安装 & \pageref{sec:software}\\
6 下单核对 & 怎样避免买错接口和重复采购 & \pageref{sec:check}\\\bottomrule
\end{tabularx}

\vfill
{\small\color{Muted}状态标记：\textbf{确定}表示资料明确支持；\textbf{推荐候选}表示规格匹配但仍需实物验证；\textbf{暂缓}表示依赖底盘、功耗或官方答疑。本文不含价格与实测性能。}

\newpagex
\section{赛题边界与总体原则}\label{sec:rules}
\tagline{先保证合规，再决定买什么。}

赛题三要求以 FPGA 构建实时控制系统，并允许视觉识别、场景理解由高性能边缘平台完成；拓展要求写作“允许使用 RK3588 等边缘 AI 平台”。指南没有出现“禁止树莓派”的文字，也没有点名保证树莓派合规。因此，树莓派可作为“等边缘 AI 平台”的合理解释，但采购前应在官方指导群取得书面确认。\cite{contest}

\begin{table}[H]\small
\caption{赛题要求与平台职责边界}\label{tab:boundary}
\begin{tabularx}{\linewidth}{@{}L{3.1cm}Y Y@{}}\toprule
\textbf{功能} & \textbf{ACG720 + 软核} & \textbf{树莓派 5}\\\midrule
摄像头与姿态显示 & 必须具备；保留 FPGA 图像通路和 IMU 数据 & 可用独立相机辅助识别\\
编码器闭环运动 & FPGA 内完成测速、PI、PWM 与安全仲裁 & 不进入实时速度环\\
测距与避障 & 距离有效性、减速、停车和基本绕障可独立运行 & 可提供目标或场景提示\\
路径点导航 & 软核保持可运行的基础导航 & 可选高层路径或语义目标\\
目标检测与任务 & 接收带时间戳的结构化结果 & 负责检测、跟踪、颜色属性、搜索\\\bottomrule
\end{tabularx}
\end{table}

\warn{购买前确认}{向官方提问：“ACG720 完成任务1--4，Raspberry Pi 5 仅用于任务5--6，通过 UART/以太网交互，是否属于允许的边缘 AI 平台？”在得到答复前，可先完成 FPGA 基础链路，树莓派专用相机可暂缓下单。}

\begin{figure}[H]\centering
\begin{tikzpicture}[node distance=7mm and 8mm]
\node[main,text width=3cm] (need) {需求确认};
\node[teal,right=of need,text width=3cm] (reuse) {核对套装已有物品};
\node[teal,right=of reuse,text width=3cm] (first) {购买首批模块};
\node[wait,below=of first,text width=3cm] (measure) {测重量、电流与空间};
\node[wait,left=of measure,text width=3cm] (motion) {底盘与电源定型};
\node[box,left=of motion,text width=3cm] (upgrade) {按瓶颈增加 AI/雷达};
\draw[flow] (need)--(reuse);\draw[flow] (reuse)--(first);\draw[flow] (first)--(measure);
\draw[flow] (measure)--(motion);\draw[flow] (motion)--(upgrade);
\end{tikzpicture}
\caption{采购顺序：先去重，再购买确定项；动力与电源必须一起确定。}
\end{figure}

\newpagex
\section{第一批：可以先购买的确定项}\label{sec:first}

\subsection{视觉、计算与传感器}

\begin{table}[H]\small
\caption{首批采购清单}\label{tab:first}
\begin{tabularx}{\linewidth}{@{}L{2.5cm}L{4.8cm}L{1cm}Y@{}}\toprule
\textbf{类别} & \textbf{推荐型号/规格} & \textbf{数量} & \textbf{购买说明}\\\midrule
FPGA 摄像头 & 小梅哥 ACM5640-V5 & 1 & DVP、SCCB、3.3V、2$\times$10 针；资料包已有初始化、DDR3、HDMI 例程\cite{cammanual,hdmimanual}\\
边缘计算 & Raspberry Pi 5，4GB & 1 & 4GB 足够开展轻量视觉；8GB 只增加内存余量，不直接提高 NPU 性能\cite{pi5}\\
散热 & Raspberry Pi Active Cooler & 1 & 持续视觉负载建议主动散热\cite{pi5}\\
桌面电源 & Raspberry Pi 官方 27W USB-C & 1 & 仅用于实验桌；车载另配稳压供电\cite{pi5}\\
存储 & 64GB microSD，A2/U3 & 1 & 正规品牌；保存系统、模型和日志\\
AI 相机 & Camera Module 3 标准彩色版 & 1 & 配 Pi 5 所需 22 针转 15 针相机排线；非 NoIR 版更适合颜色识别\cite{picam}\\
IMU & Waveshare 10 DOF IMU Sensor (D) & 1 & ICM20948 + BMP280，I\textsuperscript{2}C；磁力计远离电机\cite{imu}\\
近距测距 & Waveshare VL53L1X Distance Sensor, SKU 14929 & 4 & 前、左前、右前、后向起步；每只预留独立 SHUT 控制\cite{tof}\\\bottomrule
\end{tabularx}
\end{table}

\note{为什么两只相机}{ACM5640-V5 接 FPGA，完成赛题基础采集与显示；Camera Module 3 接树莓派，快速开展 AI。它们不是双目组合。后期若 FPGA 通过千兆以太网稳定传图，可取消树莓派独立相机链路。}

\subsection{工具、线材与显示}

\begin{table}[H]\small
\caption{先核对套装，缺少才购买}\label{tab:accessory}
\begin{tabularx}{\linewidth}{@{}L{4cm}L{1.2cm}Y@{}}\toprule
\textbf{物品} & \textbf{数量} & \textbf{核对要点}\\\midrule
高云 USB-JTAG 下载器及排线 & 1 & ACG720 套装若包含则不买；要求支持实际板载 JTAG 接口\\
ACG720 原厂电源 & 1 & 以原配优先；核对输出电压、功率、插头和极性\\
HDMI 输入显示器与线缆 & 1 & 先复用现有显示器；笔记本 HDMI 通常是输出口，不能当显示器\\
USB 读卡器 & 1 & 已有则不买；用于写入树莓派系统镜像\\
短网线、USB 数据线、端子线 & 1套 & USB 线须支持数据；高速相机信号不使用长杜邦线\\
传感器和相机固定支架 & 1套 & 保证安装刚性和视场可复现\\\bottomrule
\end{tabularx}
\end{table}

\newpagex
\section{第二批：底盘、运动与电源}\label{sec:motion}
\tagline{这部分不能拆开随意购买：电机决定驱动器、电池、保险和机械件。}

\subsection{候选运动组合}

\begin{table}[H]\small
\caption{规格明确的候选组合，不是唯一方案}\label{tab:motion}
\begin{tabularx}{\linewidth}{@{}L{2.7cm}L{4.8cm}L{1cm}Y@{}}\toprule
\textbf{物品} & \textbf{推荐候选} & \textbf{数量} & \textbf{决定条件}\\\midrule
编码器减速电机 & Pololu \#4847，25D、12V、99:1、48 CPR & 2 & 参考空载约100 rpm、堵转约5A；按车重、轮径、坡度复核\cite{motor}\\
双路电机驱动 & Cytron MDD10A Rev2.0 & 1 & 两路有刷电机、PWM+DIR、3.3V逻辑、每路10A连续；无反接保护\cite{driver}\\
编码器电平转换 & SN74LVC244A 转接电路/模块 & 1 & 编码器若以5V供电，A/B先转换到FPGA允许电平\cite{level}\\
底盘底板 & 双轮差速、预留双板与电池空间 & 1 & 不购买附带无关主控的套件\\
电机支架 & 25D 匹配支架 & 2 & 孔距、轴向和离地高度匹配\\
轮毂和车轮 & 匹配4mm输出轴 & 2套 & 轮径参与速度与里程计换算\\
万向轮/支撑轮 & 高度匹配 & 1--2 & 避免底盘前后倾斜\\\bottomrule
\end{tabularx}
\end{table}

\warn{国产替代条件}{可换国产 25D/37D 编码器电机，但商品页必须提供额定电压、减速比、输出转速、堵转电流、编码器每转计数与输出电平。缺少其中任一项，不适合直接用于闭环设计。}

\subsection{车载供电树}

\begin{figure}[H]\centering
\begin{tikzpicture}[x=1cm,y=1cm]
\node[main,text width=3.3cm] (bat) at (0,0) {成品电池组 + BMS};
\node[box,text width=3cm] (safe) at (0,-1.5) {总开关 / 急停 / 保险};
\node[teal,text width=3.5cm] (motor) at (-5,-3.1) {电机支路\\驱动器 + 双电机};
\node[teal,text width=3.5cm] (fpga) at (0,-3.1) {稳压支路\\ACG720 允许输入};
\node[teal,text width=3.5cm] (pi) at (5,-3.1) {独立 5.1V/5A 支路\\Raspberry Pi 5};
\node[label] at (0,-4.4) {三支路共地在设计点汇接；电机大电流不经过开发板；每支路单独测电压跌落和纹波};
\draw[flow] (bat)--(safe);\draw[flow] (safe)--(motor);\draw[flow] (safe)--(fpga);\draw[flow] (safe)--(pi);
\end{tikzpicture}
\caption{建议车载供电架构。具体电池串数和稳压型号必须在电机定型后确定。}\label{fig:power}
\end{figure}

\begin{table}[H]\small
\caption{电源与保护采购项}\label{tab:power}
\begin{tabularx}{\linewidth}{@{}L{4cm}L{1cm}Y@{}}\toprule
\textbf{物品} & \textbf{数量} & \textbf{下单条件}\\\midrule
成品电池组（带保护板）及匹配充电器 & 各1 & 电机额定电压、持续/峰值电流和目标续航确定后选\\
树莓派车载稳压模块 & 1 & 满足持续 5.1V/5A、瞬态和散热要求；优先选有测试报告的模块\\
ACG720 稳压模块 & 1 & 输出必须落在实物板允许范围；不要把3S“标称12V”误当恒定12V\\
保险丝、保险座、总开关、急停 & 1套 & 按支路线径、正常电流和故障电流选型\\
电源分配板、端子、粗线与接插件 & 1套 & 电机回路不能用普通杜邦线\\\bottomrule
\end{tabularx}
\end{table}

\section{感知覆盖与可选升级}\label{sec:upgrade}

\subsection{四路 ToF 的起步布局}

\begin{figure}[H]\centering
\begin{tikzpicture}[x=1cm,y=1cm]
\node[main,minimum width=5cm,minimum height=7cm] (car) at (0,0) {\begin{minipage}{4.2cm}\centering ACG720 + Pi 5\\双轮差速底盘\\[8mm]\color{white}前进方向 $\uparrow$\end{minipage}};
\node[teal] (front) at (0,4.4) {前向 ToF};
\node[teal] (fl) at (-4,3) {左前 ToF};
\node[teal] (fr) at (4,3) {右前 ToF};
\node[teal] (rear) at (0,-4.4) {后向 ToF};
\draw[flow] (car.north)--(front);\draw[flow] (car.north west)--(fl);\draw[flow] (car.north east)--(fr);\draw[flow] (car.south)--(rear);
\draw[opt] (car.west)--++(-2,0) node[left,label]{侧面/车角仍有盲区};
\draw[opt] (car.east)--++(2,0) node[right,label]{旋转前检查扫掠空间};
\end{tikzpicture}
\caption{四路 ToF 只是低速原型覆盖，不等于全向安全感知。}
\end{figure}

多只 VL53L1X 的默认地址相同，应由 FPGA/软核分别控制 SHUT，引导其逐一启动并改地址；也可拆分为多条 I\textsuperscript{2}C 总线。量产式接线前先验证上拉、电平、线长与更新时间。\cite{tof}

\subsection{暂时不买与升级触发条件}

\begin{table}[H]\small
\caption{可选升级清单}\label{tab:optional}
\begin{tabularx}{\linewidth}{@{}L{3.2cm}L{3.4cm}Y@{}}\toprule
\textbf{物品} & \textbf{现在的决定} & \textbf{何时再买}\\\midrule
Raspberry Pi AI HAT+ & 暂缓 & CPU 推理速度达不到实际任务帧率，且目标模型能转换到 Hailo 工具链时\cite{aihat}\\
二维激光雷达 & 暂缓 & 路径点导航已稳定，需要更完整的二维环境观测或定位修正时\\
车载 TFT 屏 & 按需 & 现有 HDMI 显示器完成桌面联调后，若展示要求车载显示再买 TFT4.3/5.0 V3.0（4.3/5.0 英寸）\cite{displaymanual}\\
FPGA 至 Pi 图像共享 & 后做 & 双相机方案跑通后，确有统一视角或降低硬件数量的需求\\
额外 MCU 主控板 & 不买 & 当前职责由 FPGA、软核和树莓派覆盖；新增主控会模糊 FPGA 主体地位\\
独立 DDR/Flash 扩展 & 不买 & ACG720 已板载存储，除非后续资源报告证明不足\\\bottomrule
\end{tabularx}
\end{table}

\section{软件、账号与许可证}\label{sec:software}

\begin{figure}[H]\centering
\begin{tikzpicture}[node distance=6mm]
\node[main,text width=11cm] (app) {应用层：路径点 / 任务状态机 / 目标搜索};
\node[teal,below=of app,text width=5.2cm,xshift=-3cm] (soft) {ACG720 软核\\C/C++ 固件与设备驱动};
\node[wait,below=of app,text width=5.2cm,xshift=3cm] (ai) {Raspberry Pi OS\\Picamera2 / OpenCV / 推理运行时};
\node[box,below=of soft,text width=5.2cm] (rtl) {FPGA RTL\\视频、编码器、PI、PWM、安全};
\node[box,below=of ai,text width=5.2cm] (model) {AI 模型\\转换、量化、精度/时延测试};
\draw[flow] (app)--(soft);\draw[opt] (app)--(ai);\draw[flow] (soft)--(rtl);\draw[opt] (ai)--(model);
\end{tikzpicture}
\caption{软件栈：实时控制不等待 Linux 或视觉推理。}
\end{figure}

\begin{table}[H]\small
\caption{必须或建议安装的软件}\label{tab:software}
\begin{tabularx}{\linewidth}{@{}L{3.4cm}L{2.4cm}Y@{}}\toprule
\textbf{软件} & \textbf{注册/许可} & \textbf{用途与建议}\\\midrule
Gowin EDA 商业版 & 高云账号 + License & 综合、布局布线、码流；申请免费，默认一年、绑定 MAC，官网称审核后1--3个工作日\cite{gowineda}\\
Gowin Programmer & 随工具链 & SRAM 下载、Flash 固化和断电启动验证\\
GOWIN MCU Designer & 页面提供 License 申请 & 软核 C/C++ 编译、链接与调试；先按选定软核核对支持\cite{gmd}\\
ModelSim/Questa & 合法校园/商业许可 & 教材使用其编译高云库；优先咨询学校实验室，不把安装包当授权\\
Raspberry Pi Imager / OS & 无项目账号要求 & 安装64位系统；使用当前受支持版本\cite{pi5}\\
Picamera2 / OpenCV & 开源组件 & Camera Module 3 采集、颜色处理和视觉开发\cite{picam}\\
Icarus Verilog + GTKWave & 开源可选 & 适合自写 RTL 基础仿真，不能保证替代加密 IP/DDR 模型\\
串口助手 / Wireshark & 无强制注册 & 调试 UART 和以太网；资料包内 Wireshark 版本过旧，建议用当前维护版\\\bottomrule
\end{tabularx}
\end{table}

\warn{教育版误区}{当前教育版列表出现的是 GW5AST-LV138PG484AC1/I0，而 ACG720 例程约束为 GW5AT-LV138PG484AC1/I0。名称相近不等于器件相同，建议使用已确认支持实际料号的商业版。\cite{education}}

\begin{table}[H]\small
\caption{软件准备次序}\label{tab:install}
\begin{tabularx}{\linewidth}{@{}L{1.2cm}L{4.2cm}Y@{}}\toprule
\textbf{次序} & \textbf{动作} & \textbf{完成标志}\\\midrule
1 & 注册高云账号，提交 EDA/GMD 许可 & 许可文件可导入并与目标电脑 MAC 对应\\
2 & 安装 EDA、Programmer、下载器驱动 & 能打开原厂 LED 工程并识别目标器件\\
3 & 打开 138K 现成工程但保留原副本 & 工程器件为 GW5AT-LV138PG484AC1/I0，约束无错\\
4 & 验证 SRAM 下载、Flash 固化 & 断电重启后仍执行预期程序\\
5 & 安装 Raspberry Pi OS 与相机栈 & 相机连续采集且无欠压/过热提示\\
6 & 建立 FPGA--Pi 通信协议 & 心跳、时间戳、超时停车和错误码均可验证\\\bottomrule
\end{tabularx}
\end{table}

\section{最终采购表与下单核对}\label{sec:check}

\begin{table}[H]\footnotesize
\caption{合并采购表}\label{tab:finalbom}
\begin{tabularx}{\linewidth}{@{}L{1.6cm}L{5.4cm}L{1.1cm}L{2.1cm}Y@{}}\toprule
\textbf{批次} & \textbf{物品} & \textbf{数量} & \textbf{状态} & \textbf{下单前核对}\\\midrule
已有 & ACG720、原厂电源、下载器 & 各1 & 先盘点 & 实物料号、配件、接口\\
首批 & ACM5640-V5 & 1 & 推荐确定 & 2$\times$10 DVP 版本\\
首批 & Pi 5 4GB、散热器、27W电源、64GB卡 & 各1 & 官方答疑后买 & 正品、保修、电源规格\\
首批 & Camera Module 3 + Pi 5 排线 & 各1 & 官方答疑后买 & 标准彩色版、22转15针\\
首批 & ICM20948 IMU 模块 & 1 & 推荐候选 & 3.3V逻辑、I\textsuperscript{2}C接口\\
首批 & VL53L1X SKU14929 & 4 & 推荐候选 & SHUT/INT引脚齐全\\
二批 & 编码器减速电机 + 驱动 + 电平转换 & 2+1+1 & 成套定型 & 堵转电流、编码器电平\\
二批 & 底板、支架、轮毂、轮、万向轮 & 1套 & 成套定型 & 尺寸、孔位、轴径\\
二批 & 电池、充电器、稳压、保护、线材 & 1套 & 暂缓 & 总功耗、散热、线径\\
升级 & AI HAT+、激光雷达、车载屏 & 按需 & 暂不买 & 由实测瓶颈触发\\\bottomrule
\end{tabularx}
\end{table}

\subsection{收到货后立即检查}

\begin{enumerate}
\item 拍摄包装、丝印、接口与配件；记录完整型号和版本。
\item 在限流电源或原配电源下单模块通电，先测电压，再连接信号。
\item 相机先跑原始例程；传感器逐个跑通信自检；电机先悬空低占空比试转。
\item 为所有线束标注电压、方向和端口；保存购买链接、手册、固件与驱动版本。
\item 不把一次演示成功当作验收：记录连续运行、重复启动、故障注入和温升。
\end{enumerate}

\begin{table}[H]\small
\caption{尚未取得的信息及责任动作}\label{tab:unknown}
\begin{tabularx}{\linewidth}{@{}L{4.2cm}Y@{}}\toprule
\textbf{待确认} & \textbf{应做动作}\\\midrule
树莓派是否正式属于允许的平台 & 在官方指导群 213923272 留存书面答复\\
ACG720 实物版本和接口电平 & 对照板卡丝印、原理图、通用管脚表与例程 CST\\
底盘重量、轮径、速度与坡度 & 建立机械预算后再锁电机减速比\\
整车持续、峰值和启动电流 & 分支测量后选电池、稳压器、保险和线径\\
AI 模型帧率与精度 & 在 Pi 5 CPU 实测后决定是否购买 AI HAT+\\\bottomrule
\end{tabularx}
\end{table}

\note{采购冻结条件}{正式批量下单前，至少拿到树莓派合规答复、ACG720 原理图/接口资料、电机堵转电流、整车功耗估算和机械安装图。}

\section*{参考资料}
\addcontentsline{toc}{section}{参考资料}
\begin{thebibliography}{99}\small\raggedright
\bibitem{contest} 高云半导体. 全国大学生嵌入式芯片与系统设计竞赛'2026：FPGA创新设计赛道选题指南[Z]. 用户提供文件，2026，赛题三见印刷15--17页.
\bibitem{cammanual} 小梅哥FPGA团队. 第44章：基于OV5640摄像头理论知识讲解[Z]. ACG720配套资料.
\bibitem{hdmimanual} 小梅哥FPGA团队. 第47章：OV5640摄像头采集HDMI显示屏显示[Z]. ACG720配套资料.
\bibitem{displaymanual} 小梅哥FPGA团队. 第28章：TFT显示屏驱动设计与验证[Z]. ACG720配套资料.
\bibitem{pi5} Raspberry Pi Ltd. Raspberry Pi 5 product specifications[EB/OL]. \url{https://www.raspberrypi.com/products/raspberry-pi-5/}.
\bibitem{picam} Raspberry Pi Ltd. Camera Module 3[EB/OL]. \url{https://www.raspberrypi.com/products/camera-module-3/}.
\bibitem{aihat} Raspberry Pi Ltd. AI HATs documentation[EB/OL]. \url{https://www.raspberrypi.com/documentation/accessories/ai-hat-plus.html}.
\bibitem{imu} Waveshare. 10 DOF IMU Sensor (D), ICM20948[EB/OL]. \url{https://www.waveshare.com/wiki/10_DOF_IMU_Sensor_(D)}.
\bibitem{tof} Waveshare. VL53L1X Distance Sensor, SKU 14929[EB/OL]. \url{https://www.waveshare.com/vl53l1x-distance-sensor.htm}.
\bibitem{motor} Pololu. 99:1 Metal Gearmotor 25D HP 12V with 48 CPR Encoder, \#4847[EB/OL]. \url{https://www.pololu.com/product/4847}.
\bibitem{driver} Cytron. MDD10A Dual Channel 10A Motor Driver[EB/OL]. \url{https://www.cytron.io/cytron/p-10amp-5v-30v-dc-motor-driver-2-channels}.
\bibitem{level} Texas Instruments. SN74LVC244A data sheet[EB/OL]. \url{https://www.ti.com/lit/ds/symlink/sn74lvc244a.pdf}.
\bibitem{gowineda} 高云半导体. 云源软件商业版及 License 申请[EB/OL]. \url{https://www.gowinsemi.com.cn/software/1}.
\bibitem{gmd} 高云半导体. GOWIN MCU Designer[EB/OL]. \url{https://www.gowinsemi.com.cn/software/14}.
\bibitem{education} 高云半导体. 大学计划教育版软件支持列表[EB/OL]. \url{https://www.gowinsemi.com.cn/university/software}.
\end{thebibliography}

\note{证据边界}{网页内容核验日期为2026-09-27，后续可能更新；本地教材使用 ACX720/ACG720 名称存在差异，模块兼容性仍以实物、原理图和约束文件为准。}

\end{document}
